\documentclass[11pt]{article}

\usepackage[margin=1in]{geometry}
\usepackage{amsmath, amssymb}
\usepackage{algorithm}
\usepackage{algpseudocode}
\usepackage{booktabs}
\usepackage{hyperref}

\title{Dimensionality Reduction: PCA and LDA}
\author{ML A-Z --- Part 10}
\date{}

\begin{document}
\maketitle

\section{Why Reduce Dimensionality?}

Datasets with many features ($p$ large relative to the number of
observations $n$) are hard to model and hard to visualize: models trained
on high-dimensional data are prone to overfitting (the \textbf{curse of
dimensionality}), and no plot can show more than 2--3 axes at once.
\textbf{Dimensionality reduction} techniques replace the original $p$
features with a smaller set of $k \ll p$ new features that retain as much
of the useful structure in the data as possible.

This is \textbf{feature extraction}: unlike feature \emph{selection}
(which keeps a subset of the original features), extraction builds new
features as combinations of the old ones. The two techniques covered here,
PCA and LDA, both produce new features as linear combinations of the
original ones, but they differ in what they optimize for:

\begin{itemize}
  \item \textbf{PCA} is \textbf{unsupervised} --- it ignores the target
        $y$ entirely and looks only at the variance of $X$.
  \item \textbf{LDA} is \textbf{supervised} --- it uses the class labels
        $y$ to find the directions that best separate the classes.
\end{itemize}

Both are typically applied after feature scaling (standardization), since
both rely on variance/covariance computations that are sensitive to the
scale of each feature.

\section{Principal Component Analysis (PCA)}

\subsection{Goal}

PCA looks for the directions in feature space along which the data varies
the most, under the reasoning that high-variance directions carry the most
information, while low-variance directions are closer to noise. It
projects the data onto the top $k$ such directions, called \textbf{principal
components}, discarding the rest.

\subsection{Derivation: Maximizing Variance}

Let $X \in \mathbb{R}^{n \times p}$ be the (mean-centered) data matrix. A
direction is a unit vector $\mathbf{u} \in \mathbb{R}^p$, and projecting
each observation onto it gives a new scalar feature $X\mathbf{u}$. Its
variance is

\[
  \operatorname{Var}(X\mathbf{u}) = \frac{1}{n} \mathbf{u}^\top X^\top X \, \mathbf{u}
  = \mathbf{u}^\top \Sigma \, \mathbf{u},
\]

where $\Sigma = \frac{1}{n} X^\top X$ is the ($p \times p$) covariance
matrix of the features. PCA's first component is the unit vector
$\mathbf{u}_1$ that maximizes this quantity,

\[
  \mathbf{u}_1 = \operatorname*{arg\,max}_{\lVert \mathbf{u} \rVert = 1} \mathbf{u}^\top \Sigma \mathbf{u}.
\]

Since $\Sigma$ is symmetric, it is diagonalizable with real eigenvalues:
$\Sigma \mathbf{v}_i = \lambda_i \mathbf{v}_i$, with eigenvectors
$\mathbf{v}_i$ orthonormal and eigenvalues ordered
$\lambda_1 \ge \lambda_2 \ge \dots \ge \lambda_p \ge 0$. By the Rayleigh
quotient argument, $\mathbf{u}^\top \Sigma \mathbf{u}$ is maximized (over
unit vectors) exactly at $\mathbf{u} = \mathbf{v}_1$, giving maximum value
$\lambda_1$. So:

\begin{itemize}
  \item The \textbf{principal components} are the eigenvectors of $\Sigma$,
        ordered by eigenvalue.
  \item Each eigenvalue $\lambda_i$ \emph{is} the variance captured by its
        component: $\operatorname{Var}(X\mathbf{v}_i) = \lambda_i$.
  \item Components are mutually orthogonal ($\mathbf{v}_i^\top \mathbf{v}_j = 0$
        for $i \ne j$), so they capture non-overlapping directions of
        variation --- the second component maximizes remaining variance
        subject to being orthogonal to the first, and so on.
\end{itemize}

\subsection{Algorithm}

\begin{algorithm}[h]
\caption{PCA}
\begin{algorithmic}[1]
\State Standardize $X$ (zero mean, unit variance per feature)
\State Compute the covariance matrix $\Sigma = \frac{1}{n} X^\top X$
\State Compute eigenvalues $\lambda_1 \ge \dots \ge \lambda_p$ and eigenvectors $\mathbf{v}_1, \dots, \mathbf{v}_p$ of $\Sigma$
\State Choose $k$ and form $V_k = [\mathbf{v}_1 \, \cdots \, \mathbf{v}_k]$
\State Project: $X_{\text{reduced}} = X V_k \in \mathbb{R}^{n \times k}$
\end{algorithmic}
\end{algorithm}

\subsection{Choosing \texorpdfstring{$k$}{k}: Explained Variance}

The fraction of total variance retained by the first $k$ components is

\[
  \text{explained variance ratio} = \frac{\sum_{i=1}^{k} \lambda_i}{\sum_{i=1}^{p} \lambda_i},
\]

since $\sum_i \lambda_i = \operatorname{tr}(\Sigma)$ is the total variance
in the data. A common rule of thumb is to pick the smallest $k$ that
retains e.g.\ 95\% of total variance, or to plot cumulative explained
variance against $k$ (a \textbf{scree plot}) and look for an elbow.

\section{Linear Discriminant Analysis (LDA)}

\subsection{Goal}

Where PCA asks "which directions have the most variance?", LDA asks
"which directions best separate the known classes?" It is supervised: it
uses the labels $y$ to find projections that pull different classes'
means apart while keeping each class's own spread tight.

\subsection{Scatter Matrices}

For $C$ classes, let $\boldsymbol{\mu}_c$ be the mean of class $c$ and
$\boldsymbol{\mu}$ the overall mean. Define:

\begin{itemize}
  \item \textbf{Within-class scatter} (spread inside each class, summed
        over classes):
        \[
          S_W = \sum_{c=1}^{C} \sum_{i \in c} (\mathbf{x}_i - \boldsymbol{\mu}_c)(\mathbf{x}_i - \boldsymbol{\mu}_c)^\top
        \]
  \item \textbf{Between-class scatter} (spread of the class means around
        the overall mean, weighted by class size $n_c$):
        \[
          S_B = \sum_{c=1}^{C} n_c (\boldsymbol{\mu}_c - \boldsymbol{\mu})(\boldsymbol{\mu}_c - \boldsymbol{\mu})^\top
        \]
\end{itemize}

\subsection{Derivation: Maximizing Class Separability}

LDA looks for a direction $\mathbf{u}$ that maximizes between-class
scatter relative to within-class scatter after projection, i.e.\ it
maximizes the \textbf{Fisher criterion},

\[
  J(\mathbf{u}) = \frac{\mathbf{u}^\top S_B \mathbf{u}}{\mathbf{u}^\top S_W \mathbf{u}}.
\]

Intuitively: the numerator is large when the projected class means are far
apart, and the denominator is small when each projected class is tightly
clustered around its own mean --- exactly the two properties that make
classes easy to separate. This is a generalized Rayleigh quotient, and its
maximizers are the eigenvectors of $S_W^{-1} S_B$, ordered by eigenvalue:

\[
  S_W^{-1} S_B \, \mathbf{v}_i = \lambda_i \mathbf{v}_i.
\]

These eigenvectors are the \textbf{linear discriminants}. Since $S_B$ is a
sum of $C$ rank-one-ish terms centered at a common overall mean, it has
rank at most $C - 1$, so there are at most $C - 1$ linear discriminants
with nonzero eigenvalue --- unlike PCA, LDA's number of usable components
is capped by the number of classes.

\subsection{Algorithm}

\begin{algorithm}[h]
\caption{LDA}
\begin{algorithmic}[1]
\State Standardize $X$
\State Compute per-class means $\boldsymbol{\mu}_c$ and overall mean $\boldsymbol{\mu}$
\State Compute within-class scatter $S_W$ and between-class scatter $S_B$
\State Compute eigenvalues/eigenvectors of $S_W^{-1} S_B$
\State Choose $k \le C - 1$ and form $V_k$ from the top-$k$ eigenvectors
\State Project: $X_{\text{reduced}} = X V_k$
\end{algorithmic}
\end{algorithm}

\section{PCA vs. LDA}

\begin{center}
\begin{tabular}{lll}
  \toprule
  & PCA & LDA \\
  \midrule
  Uses labels $y$? & No (unsupervised) & Yes (supervised) \\
  Optimizes & Variance of projected data & Between/within class scatter ratio \\
  Matrix diagonalized & $\Sigma$ (covariance) & $S_W^{-1} S_B$ \\
  Max useful components & $p$ (\# features) & $C - 1$ (\# classes $- 1$) \\
  Typical use & Compression, visualization, noise reduction & Preprocessing for classification \\
  \bottomrule
\end{tabular}
\end{center}

In the accompanying notebooks, both are applied to the Wine dataset (13
features, 3 cultivar classes), reduced to 2 components, and fed into a
logistic regression classifier. LDA reaches perfect test accuracy there
precisely because it optimizes directly for class separability rather than
for variance, which incidentally may or may not align with what separates
the classes.

\section{Key Takeaways}

\begin{itemize}
  \item Both PCA and LDA are linear feature-extraction methods: the new
        features are linear combinations of the old ones, found via an
        eigendecomposition.
  \item PCA is unsupervised and maximizes variance; it can find up to $p$
        components and is a good default when there is no label, or when
        the goal is compression/visualization rather than classification.
  \item LDA is supervised and maximizes class separability; it is capped
        at $C-1$ components and tends to work better as a preprocessing
        step specifically for classification, since it optimizes for the
        thing that actually matters for that task.
  \item Kernel PCA extends PCA to nonlinear structure by implicitly
        mapping data into a higher-dimensional space (via the kernel
        trick) before finding principal components --- useful when classes
        are not linearly separable in the original feature space.
\end{itemize}

\end{document}
